\section*{Discussion}

Our results show that repeated choices in a stable environment are not just repeated measurements of a stable preference. When the layout of the task is predictable, people learn to stop looking. They offload the cognitive effort of evaluating tradeoffs and rely on spatial habits. When the environment changes, attention returns, and the choices change. This has three major implications for how we understand and model decision-making.

First, choice inertia is not a static baseline; it is endogenous to attention. Standard discrete choice models often add an "inertia" dummy variable to capture why people repeat choices. But our analysis shows that time alone does not create a habit. For participants whose visual attention remained high, inertia stayed completely flat across 100 trials. Inertia only exploded for the group whose attention dropped. This means that a repeated choice is not a single phenomenon. A person might repeat a choice because they actively evaluated the attributes and still prefer it, or they might repeat it blindly because the stable layout allowed them to use a spatial shortcut. 

Second, without eye-tracking, these two behaviors look exactly the same. This is a severe case of observational equivalence. Because the decision to use a spatial heuristic is highly correlated with dropping attention, standard models that omit attention force this behavior into the error term. In our own baseline model without eye-tracking, the aggregate inertia parameter washed out and lost statistical significance. It was only by using gaze entropy to measure latent attention that we could break the tie. Eye-tracking is not just an interesting behavioral add-on; in stable environments, it is necessary to solve a fundamental endogeneity problem in the choice model.

Finally, our findings suggest that attention does more than just filter information; it dynamically constructs preferences. This provides leverage on a classic problem in neuroeconomics: do people choose an option because they look at it, or do they look at it because they already prefer it? Models like the attentional Drift Diffusion Model (aDDM) \cite{krajbich2010, krajbichRangel2011} have clearly shown that where people look predicts what they choose, but the causal arrow is often hard to pin down in static tasks. Our experimental design directly addresses this. At trial 101, the underlying economic options and their attribute levels did not change; only the spatial layout on the screen changed. There is no rational economic reason for a person's underlying preference for public transport or their willingness-to-pay for comfort to spontaneously shift at that exact moment. The only thing the manipulation forced was a return to active visual search. Because estimated willingness-to-pay shifted simultaneously with this exogenous shock to visual attention, we can isolate the causal effect. The environment forced the attentional shift, and that shift drove the change in the choice process. This supports the idea that willingness-to-pay is not a fixed number waiting to be uncovered. It is constructed in the moment, and it depends heavily on how the environment forces us to deploy our attention.

These findings connect the psychology of habits \cite{bouton2020, pool2022} with structural econometrics. For applied researchers, the lesson is practical. Long, repeated discrete choice experiments with stable layouts are not neutral measurement tools. The stable environment itself trains the participant to stop paying attention, creating the very inertia the analyst later tries to measure. If we want to capture true preferences, we have to account for the cognitive state of the person making the choice.
